\chapter{Dữ liệu, tiền xử lý và biểu diễn đặc trưng}
\label{chap:dulieu}

\section{Vai trò của dữ liệu trong bài toán}

Một mô hình học máy không trực tiếp ``hiểu'' tweet hay ảnh. Nó học quan hệ thống kê
giữa một biểu diễn số của đầu vào và nhãn mục tiêu. Vì vậy, chất lượng của dự án phụ
thuộc vào ba điều kiện trước khi lựa chọn thuật toán:

\begin{enumerate}[leftmargin=*]
  \item \textbf{đúng đối tượng:} mỗi dòng phải thực sự mô tả một cặp tweet--ảnh và
        mang đúng nhãn chính thức;
  \item \textbf{đúng ranh giới đánh giá:} thông tin của dev/test không được đi vào
        quá trình fit đặc trưng hoặc huấn luyện;
  \item \textbf{đúng biểu diễn:} văn bản và ảnh phải được chuyển thành các vector
        có ý nghĩa đối với classifier.
\end{enumerate}

Trong chương này, \textit{tiền xử lý} là các thao tác làm đầu vào nhất quán hơn,
chẳng hạn bỏ URL hoặc chuyển ảnh sang RGB. \textit{Biểu diễn đặc trưng} là bước
chuyển đầu vào thành vector số. \textit{Feature engineering} là toàn bộ quá trình
lựa chọn, tạo và kiểm soát các đặc trưng để mô hình sử dụng. Theo nghĩa đó, dự án
có hai nhánh feature engineering: TF--IDF cho văn bản và embedding CLIP cho ảnh.

\section{Bộ dữ liệu CrisisMMD v2.0}

\subsection{Nguồn, đơn vị quan sát và phạm vi}

CrisisMMD là bộ dữ liệu đa phương thức do Qatar Computing Research Institute
công bố cho nghiên cứu ứng phó thảm họa \cite{crisismmd2018}. Phiên bản sử dụng
trong dự án gồm \textbf{18.082 cặp tweet--ảnh} thuộc bảy sự kiện năm 2017:
ba cơn bão Harvey, Irma, Maria; cháy rừng California; động đất Mexico; động đất
Iraq--Iran; và lũ Sri Lanka.

Đơn vị quan sát của dự án là \textbf{một cặp tweet--ảnh}, không phải một người dùng
hay một sự kiện. Một tweet có thể xuất hiện với nhiều ảnh, do đó
\code{tweet\_id} không nhất thiết duy nhất trong tập train; cặp
\code{tweet\_id, image\_id} mới là khóa của một quan sát. Sự phân biệt này quan
trọng khi kiểm tra trùng lặp và khi căn thứ tự embedding với annotation.

\subsection{Hai nhiệm vụ dự báo}

Dự án sử dụng hai bộ nhãn chung do CrisisMMD cung cấp:

\begin{description}[leftmargin=4.2cm,style=nextline]
  \item[Informativeness] Bài đăng có chứa thông tin hữu ích cho hoạt động ứng
        phó thảm họa hay không. Đây là bài toán nhị phân gồm
        \code{informative} và \code{not\_informative}.
  \item[Humanitarian category] Nếu xét nội dung nhân đạo, bài đăng thuộc nhóm
        nhu cầu hoặc thiệt hại nào. Đây là bài toán tám lớp: người bị ảnh hưởng;
        người bị thương hoặc tử vong; người mất tích hoặc được tìm thấy; thiệt
        hại hạ tầng và tiện ích; thiệt hại phương tiện; cứu hộ, tình nguyện hoặc
        quyên góp; thông tin liên quan khác; và nội dung không mang tính nhân đạo.
\end{description}

Ngoài nhãn chung của cặp, dữ liệu còn có nhãn riêng do người gán nhãn quan sát
văn bản và ảnh. Các nhãn riêng này được dùng để chẩn đoán sự bất đồng giữa hai
phương thức, không thay thế ground truth chung khi huấn luyện mô hình vận hành.

\begin{table}[H]\centering
\caption{Các trường dữ liệu được sử dụng trong pipeline.}
\begin{tabularx}{\textwidth}{p{3.7cm}X}
\toprule
\textbf{Trường} & \textbf{Ý nghĩa và cách sử dụng}\\
\midrule
\code{event\_name} & Tên thảm họa; dùng cho EDA và kiểm tra độ ổn định theo sự kiện.\\
\code{tweet\_id, image\_id} & Khóa nhận dạng, căn thứ tự cache và kiểm tra giao nhau giữa split.\\
\code{tweet\_text, image} & Văn bản gốc và đường dẫn ảnh đầu vào.\\
\code{label} & Nhãn chung chính thức của nhiệm vụ informative.\\
\code{label\_top} & Nhãn chung chính thức của nhiệm vụ humanitarian tám lớp.\\
\code{label\_text\_*} & Nhãn chỉ dựa trên văn bản; dùng cho phân tích bất đồng.\\
\code{label\_image\_*} & Nhãn chỉ dựa trên ảnh; dùng cho phân tích bất đồng.\\
\code{multimodal\_agree} & Cờ đồng thuận category giữa người gán nhãn text và image.\\
\bottomrule
\end{tabularx}
\end{table}

\section{Hợp nhất annotation và bảo toàn split}

\subsection{Vấn đề khác partition giữa hai task}

CrisisMMD phát hành các file train/dev/test riêng cho từng nhiệm vụ. Partition
của informative và humanitarian không giống nhau. Nếu chỉ inner join hai file
cùng tên, nhiều cặp sẽ bị loại vì một mẫu có thể nằm ở train của task này nhưng
nằm ở dev của task kia. Nếu suy diễn nhãn informative chung từ nhãn
humanitarian chung, \textbf{2.649/18.082 mẫu} sẽ bị gán sai, tương đương
14,65\%.

Pipeline giải quyết theo ba bước:

\begin{enumerate}[leftmargin=*]
  \item chọn partition humanitarian làm ranh giới train/dev/test thống nhất;
  \item ghép toàn bộ annotation informative thành bảng tra cứu theo cặp
        \code{tweet\_id, image\_id};
  \item left join nhãn informative chính thức vào từng partition humanitarian
        và yêu cầu quan hệ khóa là one-to-one.
\end{enumerate}

Cách làm này giữ đủ 18.082 dòng, đồng thời hai nhiệm vụ dùng chung một ranh giới
đánh giá. Quan hệ ``khác \code{not\_humanitarian} thì informative'' chỉ đúng với
nhãn riêng từng phương thức trong dữ liệu này; dự án không dùng quan hệ đó để
tạo ground truth chung.

\subsection{Kích thước và phân bố}

\begin{table}[H]\centering
\caption{Kích thước và chất lượng cơ bản sau khi hợp nhất.}
\begin{tabular}{lrrr}
\toprule
\textbf{Chỉ số} & \textbf{Train} & \textbf{Dev} & \textbf{Test}\\
\midrule
Số cặp tweet--ảnh & 13.608 & 2.237 & 2.237\\
Giá trị thiếu & 0 & 0 & 0\\
Trùng cặp tweet--ảnh trong split & 0 & 0 & 0\\
Đường dẫn ảnh thiếu & 0 & 0 & 0\\
Số lớp humanitarian quan sát được & 8 & 8 & 8\\
Hàng qua bộ lọc exact-duplicate & 13.608 & 2.189 & 2.169\\
Hàng qua bộ lọc robustness & 13.608 & 2.078 & 2.032\\
\bottomrule
\end{tabular}
\end{table}

Trong train, informative có 8.338 mẫu và not-informative có 5.270 mẫu. Độ mất
cân bằng nghiêm trọng hơn ở task tám lớp: \code{not\_humanitarian} có 5.260 mẫu,
trong khi \code{missing\_or\_found\_people} chỉ có 24 mẫu, chênh lệch khoảng
\textbf{219:1}. Hiểu trực giác, cứ \textbf{219} bài thuộc nhóm phổ biến mới có
\textbf{1} bài thuộc nhóm hiếm. Accuracy vì vậy không đủ phản ánh khả năng phát
hiện lớp hiếm; thiết kế đánh giá phải bổ sung Macro-F1, F2 và kết quả theo từng
lớp.

\begin{insight}[Vì sao Accuracy đánh lừa khi dữ liệu lệch]
Xét một lớp hiếm có 24 mẫu trên tổng 5.284 mẫu (lớp này + nhóm phổ biến).
\emph{Mô hình A} bỏ qua hoàn toàn lớp hiếm, đoán mọi bài là nhóm phổ biến: nó
sai cả 24 bài hiếm nhưng vẫn đạt Accuracy $5.260/5.284 \approx 99{,}5\%$.
\emph{Mô hình B} bắt đúng 18/24 bài hiếm (Recall lớp hiếm $=0{,}75$) và chỉ lỡ
vài bài phổ biến, Accuracy cũng quanh $99\%$. Hai mô hình gần như cùng Accuracy
nhưng giá trị cứu hộ khác hẳn nhau: A bỏ sót toàn bộ người mất tích, B thì
không. Đây là lý do dự án khóa Macro-F1 và Recall lớp hiếm thay vì nhìn Accuracy.
\end{insight}

\section{Kiểm soát trùng lặp và rò rỉ dữ liệu}

\subsection{Rò rỉ dữ liệu là gì?}

Rò rỉ dữ liệu xảy ra khi quá trình huấn luyện tiếp cận thông tin mà trong sử dụng
thực tế chỉ xuất hiện ở tương lai hoặc thuộc tập đánh giá. Khi đó, chỉ số test có
thể cao nhưng không phản ánh khả năng tổng quát hóa. Với mạng xã hội, chỉ kiểm
tra ID là chưa đủ: hai bài đăng khác ID vẫn có thể dùng cùng một ảnh hoặc cùng
nội dung sau khi bỏ URL và mention.

Dự án kiểm tra bốn mức:

\begin{enumerate}[leftmargin=*]
  \item giao nhau của \code{tweet\_id}, \code{image\_id} và cặp ID;
  \item trùng ảnh chính xác bằng SHA-256;
  \item trùng văn bản sau bước làm sạch;
  \item gần trùng bằng perceptual hash cho ảnh và độ tương tự ngữ nghĩa/từ vựng
        cho văn bản.
\end{enumerate}

Không có ID hoặc cặp ID giao nhau giữa các split. Tuy nhiên, SHA-256 phát hiện
43 hash ảnh chung train--dev, 46 hash chung train--test và 21 hash chung
dev--test. Có một văn bản sau làm sạch trùng giữa train với dev và một văn bản
trùng giữa train với test. Trên toàn bộ corpus, 18.082 dòng chỉ tạo ra 17.777
hash ảnh duy nhất, tức 305 lần lặp ngoài ảnh đại diện đầu tiên.

\subsection{Hai loại ``vân tay'' ảnh: SHA-256 và perceptual hash}

Để phát hiện ảnh trùng mà không phải so từng cặp ảnh (rất tốn kém), dự án dùng
hai loại hàm băm (hash) bổ sung cho nhau. Một \emph{hàm băm} biến một đầu vào
thành một chuỗi bit ngắn cố định, gọi là ``vân tay''.

\begin{description}[leftmargin=3.6cm,style=nextline]
  \item[SHA-256] Đọc \emph{toàn bộ byte} của file ảnh và tạo một vân tay 256 bit.
        Cùng một file (từng byte giống nhau) luôn cho cùng hash; chỉ cần đổi một
        byte, hash đã khác hẳn. SHA-256 vì vậy phát hiện ảnh \emph{giống hệt},
        nhưng \emph{không} biết hai ảnh ``trông giống nhau''.
  \item[Perceptual hash (pHash)] Rút đặc trưng thị giác tần số thấp (bố cục
        sáng/tối tổng thể) rồi tạo một vân tay 64 bit. Ảnh bị resize hoặc crop
        nhẹ vẫn cho hash \emph{gần} nhau. Mức ``gần'' đo bằng \emph{khoảng cách
        Hamming} -- số bit khác nhau giữa hai hash.
\end{description}

\begin{figure}[H]
\centering
\begin{tikzpicture}[
  font=\small,
  imgA/.style={draw,minimum width=1.1cm,minimum height=1.1cm,fill=blue!12},
  imgB/.style={draw,minimum width=1.1cm,minimum height=1.1cm,fill=blue!25},
  hash/.style={draw=HUSTBlue,rounded corners=1.5pt,fill=boxgray,font=\ttfamily\scriptsize,
    minimum height=0.7cm},
  lab/.style={font=\scriptsize\itshape}
]
% Row 1: exact duplicate via SHA-256
\node[imgA] (a1) at (0,1.6) {};
\node[imgA] (a2) at (1.4,1.6) {};
\node[hash,right=0.5cm of a2] (ha) {9f3a\dots e1 = 9f3a\dots e1};
\node[lab,right=0.3cm of ha] {SHA-256 khớp $\Rightarrow$ giống hệt};
\node[lab,left=0.2cm of a1] {Ảnh, copy};
% Row 2: near-duplicate via pHash
\node[imgA] (b1) at (0,0) {};
\node[imgB] (b2) at (1.4,0) {};
\node[hash,right=0.5cm of b2] (hb) {1100\dots01 vs 1100\dots11};
\node[lab,right=0.3cm of hb] {Hamming $=2 \Rightarrow$ gần trùng};
\node[lab,left=0.2cm of b1] {Ảnh, crop};
\end{tikzpicture}
\caption{Hai lớp phát hiện trùng lặp. SHA-256 chỉ bắt ảnh giống hệt từng byte;
pHash bắt ảnh gần giống qua khoảng cách Hamming nhỏ.}
\label{fig:hash-explainer}
\end{figure}

\subsection{Hai mặt nạ đánh giá}

\begin{description}[leftmargin=3.8cm,style=nextline]
  \item[Canonical mask] Loại một dòng dev/test nếu ảnh SHA-256 hoặc văn bản đã
        làm sạch của nó đã xuất hiện ở split trước. Đây là bộ lọc tự động có
        tiêu chí xác định và được dùng trong quy trình chọn mô hình ban đầu.
  \item[Robustness mask] Bắt đầu từ canonical mask, sau đó loại thêm các cặp gần
        trùng đã được duyệt theo từng cặp. Ảnh ứng viên dùng pHash DCT 64-bit
        với khoảng cách Hamming không quá 4; văn bản dùng kết hợp cosine của
        CLIP text embedding và TF--IDF ký tự/từ. Mặt nạ này dùng để kiểm tra độ
        bền của kết luận, không tự động coi mọi ứng viên là bản sao.
\end{description}

Sau duyệt cặp, robustness mask loại thêm 111 dòng dev và 137 dòng test so với
canonical mask. Việc báo cáo cả kết quả canonical và robustness giúp phân biệt
hai câu hỏi: mô hình hoạt động thế nào trên split đã làm sạch bản sao chính xác,
và kết luận còn giữ được hay không khi loại cả những bản sao gần.

\section{Tiền xử lý và feature engineering văn bản}

\subsection{Làm sạch tweet}

Tweet chứa nhiều token phục vụ nền tảng hơn là nội dung: URL rút gọn, mention,
HTML entity và dấu câu. Hàm làm sạch thực hiện lần lượt:

\begin{enumerate}[leftmargin=*]
  \item chuyển chữ thường;
  \item bỏ HTML entity, URL và \code{@mention};
  \item bỏ ký hiệu \code{\#} nhưng giữ từ khóa hashtag;
  \item chỉ giữ chữ cái Latin, chữ số và khoảng trắng;
  \item gộp nhiều khoảng trắng liên tiếp.
\end{enumerate}

Ví dụ:

\begin{center}
\begin{tabularx}{0.94\textwidth}{p{2cm}X}
\toprule
\textbf{Trước} & \code{URGENT:\ Flood\ waters\ rising!\ @Police\ https://t.co/x\ \#Harvey}\\
\textbf{Sau} & \code{urgent\ flood\ waters\ rising\ harvey}\\
\bottomrule
\end{tabularx}
\end{center}

Giữ nội dung hashtag là quyết định có chủ đích vì tên thảm họa hoặc chủ đề cứu
trợ thường nằm trong hashtag. Ngược lại, bỏ URL và mention giúp giảm số token
chỉ xuất hiện một lần.

Một câu hỏi quan trọng: \emph{còn loại rác nào không nằm trong các bước làm sạch
ở trên thì sao?} Một bộ quy tắc regex \textbf{không thể} liệt kê mọi dạng nhiễu
của ngôn ngữ mạng xã hội. Vì vậy báo cáo phân loại rõ ba mức xử lý ở
Bảng~\ref{tab:cleaning-audit} thay vì tuyên bố dữ liệu đã ``sạch hoàn toàn''.

\begin{table}[H]\centering\small
\caption{Ba mức xử lý nhiễu trong văn bản tweet.}
\label{tab:cleaning-audit}
\begin{tabularx}{\textwidth}{p{3.6cm}X}
\toprule
\textbf{Mức xử lý} & \textbf{Loại nhiễu}\\
\midrule
Đã xử lý trực tiếp & URL, \code{@mention}, HTML entity, ký hiệu \code{\#},
ký tự không phải chữ/số.\\
Giảm tác động (không loại hẳn) & Token quá hiếm bị cắt bởi giới hạn 2.000 đặc
trưng và \code{min\_df}; stopword tiếng Anh bị bỏ trọng số.\\
Chưa xử lý tốt & Emoji, mỉa mai và phủ định, lỗi chính tả, văn bản đa ngôn ngữ,
tweet chỉ có URL (rỗng sau làm sạch), spam lặp từ.\\
\bottomrule
\end{tabularx}
\end{table}

Một audit định lượng trên dữ liệu thật (lưu ở \code{text\_cleaning\_audit.json})
cho thấy: trung bình \textbf{40\%} ký tự mỗi tweet bị loại khi làm sạch -- phần
lớn là URL và mention; \textbf{không có} tweet nào trở thành chuỗi rỗng sau làm
sạch (mọi tweet vẫn còn nội dung), nên cơ chế ``tweet chỉ có URL thành rỗng'' là
một rủi ro lý thuyết chứ không xảy ra trong corpus này; và khi chuyển sang
dev/test, tỷ lệ token nằm ngoài từ vựng train (OOV) là \textbf{7,6\%} và
\textbf{7,8\%} -- đủ thấp để TF--IDF vẫn phủ phần lớn nội dung.

Phần rác còn lại chủ yếu mang tính ngữ nghĩa: một tweet mỉa mai
(``great, another flood'') có thể bị TF--IDF hiểu sai. Đây là lý do các ca như
vậy được đẩy sang Manual Review thay vì tin tuyệt đối vào dự báo (xem
Mục~\ref{sec:scenarios}). Việc loại emoji, dấu tiếng nước ngoài và không xử lý
mỉa mai là một đánh đổi có chủ đích để giữ pipeline đơn giản, kiểm tra được và
phù hợp corpus chủ yếu tiếng Anh.

\subsection{TF--IDF}

Classifier không nhận chuỗi ký tự trực tiếp. TF--IDF chuyển mỗi văn bản thành
vector thưa, trong đó một chiều tương ứng với một từ hoặc cặp từ. Trọng số của
term \(t\) trong document \(d\) được mô tả bởi:

\[
\operatorname{tfidf}(t,d)
=\operatorname{tf}(t,d)\times
\log\left(\frac{N+1}{\operatorname{df}(t)+1}\right)+1,
\]

trong đó \(\operatorname{tf}(t,d)\) phản ánh số lần term xuất hiện,
\(\operatorname{df}(t)\) là số văn bản chứa term và \(N\) là số văn bản train.
Trực giác là một từ xuất hiện nhiều trong một tweet nhưng hiếm trên toàn corpus
sẽ mang nhiều thông tin phân biệt hơn từ phổ biến ở mọi tweet.

Các cấu hình thực tế của vectorizer là:

\begin{itemize}[leftmargin=*]
  \item tối đa 2.000 đặc trưng có thứ tự theo thống kê của tập train;
  \item unigram và bigram, tức một từ và hai từ liên tiếp;
  \item stopword tiếng Anh của scikit-learn;
  \item chuẩn hóa \(L_2\) mặc định cho từng vector.
\end{itemize}

Bigram giúp phân biệt các cụm như \code{power\ outage}, \code{missing\ people}
hoặc \code{need\ shelter}, trong khi unigram giữ độ phủ khi câu ngắn. Giới hạn
2.000 chiều giảm chi phí và nhiễu so với giữ toàn bộ từ vựng, nhưng có thể bỏ
sót term hiếm. Vectorizer \textbf{chỉ fit trên train}; dev và test chỉ gọi
\code{transform}. Nếu fit trên toàn corpus, tần suất tài liệu của dev/test sẽ
đi vào IDF và gây rò rỉ nhẹ nhưng có thật.

\paragraph{Ví dụ tính tay.} Xét ba tweet đồ chơi đã làm sạch (\(N=3\)):
\(d_1=\) ``flood rescue now'', \(d_2=\) ``flood power outage'',
\(d_3=\) ``donation drive''. Từ \emph{flood} xuất hiện ở 2 văn bản
(\(\operatorname{df}=2\)), còn \emph{rescue} chỉ ở 1 (\(\operatorname{df}=1\)).
Theo công thức (dạng smooth của scikit-learn):
\[
\operatorname{idf}(\text{flood})=\ln\frac{1+3}{1+2}+1=1{,}288,\qquad
\operatorname{idf}(\text{rescue})=\ln\frac{1+3}{1+1}+1=1{,}693.
\]
Mỗi từ xuất hiện một lần trong \(d_1\) nên \(\operatorname{tf}=1\), suy ra
\(\operatorname{tfidf}(\text{flood},d_1)=1{,}288\) và
\(\operatorname{tfidf}(\text{rescue},d_1)=1{,}693\). Như vậy \emph{rescue}
(hiếm hơn) được trọng số cao hơn \emph{flood} (phổ biến hơn) -- đúng trực giác:
từ hiếm phân biệt tốt hơn. Sau chuẩn hóa \(L_2\) cho vector \(d_1\) (ba từ
flood/rescue/now), trọng số flood là \(1{,}288/2{,}719=0{,}47\) và rescue là
\(1{,}693/2{,}719=0{,}62\). Ma trận document--term tương ứng (chưa chuẩn hóa)
là một ma trận thưa:

\begin{table}[H]\centering\small
\caption{Ma trận document--term TF--IDF của ba tweet đồ chơi (ô trống $=0$).}
\begin{tabular}{lccccccc}
\toprule
& \textbf{flood} & \textbf{rescue} & \textbf{now} & \textbf{power}
& \textbf{outage} & \textbf{donation} & \textbf{drive}\\
\midrule
\(d_1\) & 1,29 & 1,69 & 1,69 & & & &\\
\(d_2\) & 1,29 & & & 1,69 & 1,69 & &\\
\(d_3\) & & & & & & 1,69 & 1,69\\
\bottomrule
\end{tabular}
\end{table}

Điểm yếu cần nhớ: TF--IDF chỉ đếm từ và cặp từ, \emph{không} hiểu thứ tự xa,
phủ định (``no rescue needed'') hay mỉa mai, và xem các từ là độc lập. Đây là lý
do nhánh văn bản cần được bổ sung bằng ảnh và bằng Manual Review cho ca mơ hồ.

\section{Tiền xử lý và feature engineering ảnh}

\subsection{Từ pixel đến vector đặc trưng}

Ảnh có hàng trăm nghìn giá trị pixel, quá nhiều để các classifier cổ điển học
trực tiếp trên 13.608 mẫu train. Dự án sử dụng CLIP ViT-B/32
\cite{clip2021} như một \textbf{bộ trích đặc trưng tiền huấn luyện}. Các khái
niệm cần phân biệt:

\begin{description}[leftmargin=4.4cm,style=nextline]
  \item[Tiền huấn luyện] Mô hình CLIP đã được tác giả gốc học trước trên lượng
        lớn cặp ảnh--văn bản, trước khi dự án này sử dụng nó.
  \item[Transfer learning] Tri thức hình ảnh đã học được chuyển sang bài toán
        CrisisMMD thay vì huấn luyện mạng sâu từ đầu.
  \item[Frozen feature extractor] Trọng số CLIP được giữ nguyên; dự án chỉ chạy
        forward pass để lấy vector. Không có gradient và không cập nhật CLIP.
  \item[Fine-tuning] Là cập nhật một phần hoặc toàn bộ trọng số mô hình tiền
        huấn luyện bằng dữ liệu đích. \textbf{Dự án không fine-tune CLIP}.
\end{description}

Trong ViT-B/32, ảnh \(224\times224\) được chia thành lưới \(7\times7\) gồm
\(49\) patch kích thước \(32\times32\). Mỗi patch được biến thành một token và
đi qua Transformer để tổng hợp quan hệ giữa các vùng ảnh, cuối cùng thu về một
vector \(512\) chiều (Hình~\ref{fig:vit-clip}a). CLIP được học bằng
\emph{mục tiêu tương phản} (contrastive): trong một lô ảnh--caption, vector của
mỗi ảnh được kéo gần caption đúng của nó và đẩy xa các caption còn lại
(Hình~\ref{fig:vit-clip}b). Nhờ vậy embedding chứa thông tin ngữ nghĩa cao hơn
màu sắc hay histogram đơn giản.

\begin{figure}[H]
\centering
\begin{subfigure}{0.54\textwidth}
\centering
\begin{tikzpicture}[font=\scriptsize,node distance=4mm]
% 7x7 patch grid
\foreach \i in {0,...,6}{
  \foreach \j in {0,...,6}{
    \fill[blue!\the\numexpr10+\i*6+\j\relax] (\i*0.28,\j*0.28) rectangle ++(0.26,0.26);
  }
}
\draw[HUSTBlue,thick] (0,0) rectangle (1.96,1.96);
\node[below=1mm] at (0.98,0) {Ảnh $224{\times}224$ $\to$ $7{\times}7$ patch};
\node[draw=HUSTBlue,rounded corners=2pt,fill=boxgray,right=1.1cm of current bounding box.east,
  minimum height=0.8cm,align=center] (tf) at (3.6,1) {Transformer\\(49 token)};
\node[draw=HUSTBlue,rounded corners=2pt,fill=boxgray,below=5mm of tf,align=center] (vec)
  {Vector $512$-D};
\draw[-{Latex[length=2mm]},thick,HUSTBlue] (1.96,1) -- (tf.west);
\draw[-{Latex[length=2mm]},thick,HUSTBlue] (tf) -- (vec);
\end{tikzpicture}
\caption{Patch $\to$ token $\to$ Transformer $\to$ embedding.}
\end{subfigure}\hfill
\begin{subfigure}{0.42\textwidth}
\centering
\begin{tikzpicture}[font=\scriptsize]
\foreach \i in {0,1,2}{
  \foreach \j in {0,1,2}{
    \ifnum\i=\j
      \fill[HUSTBlue!70] (\j*0.7,-\i*0.7) rectangle ++(0.6,0.6);
    \else
      \fill[black!8] (\j*0.7,-\i*0.7) rectangle ++(0.6,0.6);
    \fi
  }
}
\node[left] at (0,-0.3) {ảnh 1};
\node[left] at (0,-1.0) {ảnh 2};
\node[left] at (0,-1.7) {ảnh 3};
\node[above] at (0.3,0.6) {cap.1};
\node[above] at (1.0,0.6) {cap.2};
\node[above] at (1.7,0.6) {cap.3};
\end{tikzpicture}
\caption{Ô chéo (cặp đúng) được kéo gần, ô ngoài bị đẩy xa.}
\end{subfigure}
\caption{Cách CLIP biến ảnh thành vector và cách nó được huấn luyện tương phản.}
\label{fig:vit-clip}
\end{figure}

Một ví dụ trực giác về lân cận: trong không gian CLIP, vector của một ảnh
\emph{đường ngập} thường gần vector của một ảnh ngập khác hơn là gần ảnh
\emph{thùng quyên góp}. Cần nhấn mạnh: embedding \textbf{không phải} là xác suất
hay nhãn; nó chỉ là tọa độ ngữ nghĩa. Chính classifier phía sau (Chương~\ref{chap:mohinh})
mới học ranh giới giữa các lớp CrisisMMD trên không gian đó. Và như đã nêu, CLIP
luôn được giữ \textbf{frozen, không fine-tune}.

\subsection{Quy trình trích xuất thực tế}

Mỗi ảnh được:

\begin{enumerate}[leftmargin=*]
  \item mở bằng Pillow, chuyển về ba kênh RGB và resize \(224\times224\);
  \item đưa qua bộ xử lý chính thức của CLIP và forward pass ở chế độ
        \code{eval} trong \code{torch.no\_grad()};
  \item lấy projected image embedding 512 chiều;
  \item chuẩn hóa \(L_2\):
        \(\hat{\mathbf{x}}=\mathbf{x}/\lVert\mathbf{x}\rVert_2\);
  \item lưu cache theo đúng thứ tự \code{tweet\_id, image\_id}.
\end{enumerate}

Chuẩn hóa \(L_2\) đưa mọi vector lên mặt cầu đơn vị, giúp khoảng cách và tích vô
hướng phản ánh hướng ngữ nghĩa thay vì độ lớn tùy ý. Cache có shape
\((13.608,512)\), \((2.237,512)\), \((2.237,512)\) cho train/dev/test. Audit
cho thấy norm trung bình bằng 1, độ lệch chuẩn toàn ma trận khoảng 0,04418,
không có vector zero và metadata khớp thứ tự annotation.

\begin{figure}[H]\centering
\includegraphics[width=0.92\textwidth]{image_gallery_seeded.png}
\caption{Hai ảnh train cho mỗi lớp trong tám lớp humanitarian, lấy mẫu xác định
bằng seed 42 (\emph{không} chọn tay theo ảnh ``đẹp''), nên gallery không bị thiên
lệch chủ quan.}
\end{figure}

\subsection{Lý do giữ CLIP cố định}

Giữ CLIP cố định phù hợp với phạm vi dự án vì:

\begin{itemize}[leftmargin=*]
  \item mục tiêu học phần là so sánh các thuật toán phân lớp và xây dựng DSS,
        không phải thiết kế mạng sâu;
  \item 13.608 mẫu train và các lớp hiếm rất nhỏ làm fine-tuning toàn mạng có
        nguy cơ overfit;
  \item embedding được tính một lần, sau đó cùng một không gian đặc trưng được
        dùng để so sánh công bằng sáu classifier;
  \item chi phí huấn luyện và triển khai thấp hơn, kết quả dễ tái lập hơn.
\end{itemize}

Đánh đổi là đặc trưng không được điều chỉnh riêng cho ảnh thảm họa. Kết quả của
nhánh ảnh vì vậy phải được hiểu là ``classifier trên đặc trưng CLIP cố định'',
không phải thành tích của một CLIP đã được huấn luyện cho CrisisMMD.

\section{Tổng hợp đặc trưng đầu vào mô hình}

\begin{table}[H]\centering
\caption{Đặc trưng được tạo và nơi chúng được sử dụng.}
\begin{tabularx}{\textwidth}{p{2.4cm}p{3.3cm}p{2.4cm}X}
\toprule
\textbf{Nguồn} & \textbf{Biểu diễn} & \textbf{Kích thước} & \textbf{Vai trò}\\
\midrule
Tweet & TF--IDF word 1--2 gram & 2.000, vector thưa &
Đầu vào sáu classifier văn bản cho hai task.\\
Ảnh & CLIP ViT-B/32 frozen embedding & 512, vector đặc &
Đầu vào sáu classifier ảnh cho hai task.\\
Xác suất mô hình & Vector xác suất đã căn lớp & 2 hoặc 8 &
Đầu vào Late Fusion và conflict score.\\
Từ khóa nghiệp vụ & Cờ xuất hiện trong tweet đã làm sạch & một số cờ nhị phân &
Tín hiệu bổ sung cho Risk Score, không dùng để thay đổi ground truth.\\
\bottomrule
\end{tabularx}
\end{table}

Như vậy, feature engineering của dự án không phải là một bước trang trí trước
mô hình. Nó quyết định loại thông tin mà classifier có thể nhìn thấy, chi phí
tính toán, nguy cơ rò rỉ và khả năng giải thích. Chương tiếp theo sử dụng chính
các biểu diễn này để phân tích cấu trúc dữ liệu trước khi huấn luyện.
